\begin{textsample}{POS Dim 4 – persona\_gpt – Score 24.00 – t886\_gpt.txt}  \label{ex:f4_pos_030}
Standing on the \textbf{shores} of Svalbard, the Arctic doesn't feel distant or abstract.
It feels immediate, fragile, and under pressure.
Jennifer Morgan, Greenpeace International ’s Executive Director, travelled there to \textbf{witness} how quickly the Arctic is changing as the climate crisis accelerates.
The \textbf{ice} is melting, \textbf{temperatures} are rising, and the impacts are no longer hidden in remote corners – they are written across the \textbf{landscape} and the \textbf{lives} of the \textbf{animals} that depend on it.
Along the \textbf{coastline}, one of the most visible scars is the trash washing up on the \textbf{beaches}.
Most of it comes from industrial \textbf{fisheries}.
What might look like a remote, pristine \textbf{shore} is in reality littered with ropes, \textbf{nets}, and plastic debris.
This pollution is not just an eyesore ; it is a direct \textbf{threat} to \textbf{life} in the Arctic.
Reindeer, polar bears, seabirds and other \textbf{wildlife} are being harmed by this constant stream of waste. \textbf{Animals} become entangled in \textbf{fishing} gear or ingest plastic, mistaking it for food.
These encounters can be deadly, adding another \textbf{layer} of pressure to \textbf{species} already struggling with rapidly changing \textbf{sea} \textbf{ice} and \textbf{habitat} loss driven by global heating.
To confront this, a \textbf{beach} clean-up was organised under Greenpeace ’s ‘ Protect What You Love ’ campaign.
The clean-up did more than remove trash from one stretch of \textbf{coastline} ; it helped expose the scale of the problem and connect it to the wider issue of destructive \textbf{fishing} practices, particularly \textbf{bottom} trawling.
This \textbf{method} of \textbf{fishing} drags heavy gear across the seafloor, damaging fragile marine \textbf{ecosystems} and contributing to the broader crisis facing the Arctic Ocean.
There has been a step forward.
Through Greenpeace ’s efforts, a voluntary agreement has been brokered within the \textbf{fishing} industry to avoid expanding cod \textbf{fishing} into untouched \textbf{areas} of the \textbf{northern} Barents Sea.
This is a recognition from parts of the industry that some \textbf{areas} must remain off-limits if we are to give the Arctic a chance.
But voluntary measures are not enough to match the speed and severity of the changes underway.
Jennifer Morgan is calling on Norway to turn this recognition into lasting \textbf{protection} by establishing strong, legal \textbf{safeguards} for marine \textbf{areas} in the Barents Sea.
Without binding rules, the Arctic ’s remaining relatively untouched \textbf{waters} remain vulnerable to future industrial expansion.
What is happening in Svalbard is a stark reminder that the Arctic is on the frontline of the climate and nature crisis.
Protecting it means confronting both the pollution and the destructive practices that are pushing this unique \textbf{region} – and the \textbf{wildlife} that calls it home – to the \textbf{brink}.

% matched lemmas: animal, area, beach, bottom, brink, coastline, ecosystem, fishery, fishing, habitat, ice, landscape, layer, life, method, net, northern, protection, region, safeguard, sea, shore, specie, temperature, threat, water, wildlife, witness
\end{textsample}
